\documentclass[12pt]{article} 
\renewcommand{\thempfootnote}{\arabic{mpfootnote}}

\usepackage[utf8]{inputenc}

\usepackage[T1]{fontenc}
\usepackage[ttdefault=true]{AnonymousPro}

\usepackage{multicol}
\setlength{\columnsep}{1cm}

\usepackage{abc}
\usepackage{graphicx}
\usepackage{musicography}

\usepackage[margin=0.5in]{geometry} 
\geometry{letterpaper}

\usepackage{fancyhdr} 
\pagestyle{fancy} % options: empty , plain , fancy
\renewcommand{\headrulewidth}{0pt} % customise the layout...
\lhead{}\chead{}\rhead{}
\lfoot{}\cfoot{\thepage}\rfoot{}

\begin{document}

\thispagestyle{empty}

\begin{abc}[name=catcameback,options={-O= +c -F single.fmt}]
X:1
T:And The Cat Came Back
S:Greg Canote, "Dogs in the Dishes"
C:from Doc Roberts
M:C|
L:1/8   
K:Gmaj
P:A
|: ef | "G"gage d2B2 | gagd e2 d2 |  g2-ga gdef | "D"geae "G"g2 :|
P:B
|:Bc | "G"d2 g2-ge g2 | b4-b2 _b2 | =b2 b2-bage | d2-d_d =d_d=d_d | 
=d_d=dd edBA | "D"B2-Bd BAGA | "D"BAdA BG A2 | "G" G2-GG G2 :|
P:C
|: (3DEF | "G"GFGD BcdB | "C"c2Bc dcBA | "G"GFGG BcdB | "D"ABAF DFAF | 
"G"GFGA BcdB | "C"cBcd e2 d2 | "D"edef gedA | BdAF "G"G2 :|
\end{abc}
\vskip 0.5in
\noindent\makebox[\textwidth][c]{%
  \begin{minipage}{7.0in}
    \setlength{\parindent}{1.5em}

    This transcription of ``And The Cat Came Back'' is from the
    playing of Greg Canote on the Canote Brothers' album ``Dogs in the
    Dishes''. The tune comes from a 1927 recording of Doc
    Roberts\footnote{\tt https://www.slippery-hill.com/content/and-cat-came-back},
    although on that recording he mostly plays the B part once each
    time around, not twice.

    Harry S. Miller published a song named ``The Cat Came Back'' in
    1893 which Fiddlin' John Carson recorded a version of in 1924\footnote{\tt https://youtu.be/3SA5O0acXpk}.
    That song has two parts, one of which is similar to the ``A'' part
    of this tune, so it's reasonable to assume that Doc Roberts had
    heard the popular song and was probably referencing it in his
    version, but if there's any similarity between the other parts of
    Roberts' tune and Miller's song, I'm not hearing it. The song
    itself has mutated quite a bit over the years as well; the melody
    used on The Muppet Show is completely different from the melody I
    learned at summer camp, and both are significantly different from
    the original sheet music.

    R. P. Christeson, in his 1973 \textit{The Old-Time Fiddler's
      Repertory}, wrote that this tune was ``[p]layed by a few
    Missouri fiddlers in the depression days of the early 1930's, but
    it is seldom heard any more.'' It's still not a particularly
    common tune, at least in my experience---I don't think I've played
    it at a jam without being the person who called it. Maybe
    featuring it here will help change that.

  \end{minipage}
}


\end{document}
